% Fragmento ES para inserción, no documento autónomo.
% Destino: VII/manuscrito/main.tex, después de relaciones_termica.tex.
% La sección recibe las salidas y pruebas de normalizacion_conjunta.tex,
% composicion.tex y relaciones_termica.tex; no altera sus generadores.
\section{Composición de horizonte y restitución térmico-geométrica}
\label{ins29:vii:horizonte}

La composición de las constantes permite distinguir qué conserva cada
lectura física del mismo estado generado. Recibimos las secciones de
acción y longitud, el par angular y la velocidad constitutiva ya
construidos desde APP--TRIT--TPK. Abreviamos \(\eta=\eta_{\rm ret}\)
y denotamos por \(s_0\) la sección de acción cuya expresión en la coordenatización
SI es \(10^{-34}\,\mathrm{J\,s}\). Así,
\[
 \hbar=\eta s_0,\qquad
 L_\alpha=r_N\alpha^{16}L_*,\quad r_N=\frac{5759}{23040},\qquad
 G=\frac{c^3L_\alpha^2}{\eta s_0}.
\]
El retorno inscrito determina \(r_N\); el normando local determina la
potencia \(16\); la sección \(L_*\) conserva la coordenatización longitudinal.
La función \(\eta=H_5-90\mathcal C_\pi/\pi\), con sus coeficientes y
dependencias regionales, es la obtenida en
\eqref{md:compos:seccion-accion}. Las sustituciones siguientes reciben esas
salidas y conservan las unidades y los dominios de sus lectores.

\subsection{Área fija, masa fija y conservación del producto térmico}

Para el contraste de horizonte de Einstein, la expresión de
Bekenstein--Hawking es
\(S_{\rm BH}=k_Bc^3\mathcal A/(4G\hbar)\).
El coeficiente \(1/4\) pertenece aquí a esta ley de contraste.
En un horizonte de Schwarzschild, estacionario, sin carga ni rotación,
\(R_S=2GM/c^2\), \(\mathcal A=4\pi R_S^2\) y
\(k_BT_H=\hbar c/(4\pi R_S)\).

\begin{proposition}[Dos restricciones de la composición de horizonte]
\label{ins29:vii:area_masa}
En la realización indicada, las salidas estructurales dan
\begin{align*}
 S_{\mathcal A}&=\frac{k_B\mathcal A}{4L_\alpha^2},&
 T_{\mathcal A}&=\frac{\eta s_0c}{2k_B\sqrt{\pi\mathcal A}},\\
 S_M&=\frac{4\pi k_Bc^2L_\alpha^2M^2}{\eta^2s_0^2},&
 T_M&=\frac{\eta^2s_0^2}{8\pi k_BL_\alpha^2M},
 \qquad S_MT_M=\frac{Mc^2}{2}.
\end{align*}
En la comparación entre secciones que conserva \(c,L_\alpha,k_B\),
la lectura a área fija mantiene \(S_{\mathcal A}\) y hace
\(T_{\mathcal A}\) proporcional a \(\eta\); la lectura a masa fija
da \(S_M\propto\eta^{-2}\) y \(T_M\propto\eta^2\).
\end{proposition}
\begin{proof}
La identidad \(G\hbar=c^3L_\alpha^2\) da inmediatamente la primera
entropía. Sustituir \(R_S=\sqrt{\mathcal A/(4\pi)}\) en la temperatura
da la primera temperatura. A masa fija,
\(R_S=2cL_\alpha^2M/(\eta s_0)\); su área y su inverso producen las
otras dos fórmulas. Multiplicarlas cancela los factores comunes y deja
\(Mc^2/2\).
\end{proof}

La cancelación de la acción en la entropía a área fija expresa una
compatibilidad entre sus lecturas gravitatoria y longitudinal. A masa
fija, la misma sección de acción interviene en la geometría del horizonte,
y su dependencia permanece en la entropía. El balance debe conservar la
restricción elegida durante toda la comparación.
El parámetro radial \(R\) de la realización anterior
\(\hbar c/(2\pi R)\) coincide en este contraste con \(2R_S\):
ambas expresiones representan entonces la misma temperatura. Esta
identificación conserva el factor dos entre radio geométrico y escala
de aceleración.

\subsection{Separación angular y restitución de la sección de acción}

La realización modular de frontera conserva dieciocho ocupaciones
tríticas en cada sector de incidencia. Escribimos \(X=N_{90}\),
\(Y=N_{120}\), \(N=X+Y\) y \(M=X-Y\). Su separación declarada es
\[
 \Delta_{\rm mod}=
 \frac{(A_{\deg}-C^*_{\deg})\pi/180}{9\cdot90}
 \left(1-\frac7{12}\frac\pi{729}\right),\qquad
 d=30\Delta_{\rm mod},\quad f_\pi=1-\frac{7\pi}{8748}.
\]
La coordenada \(d\) conserva tanto la conversión angular \(\pi/180\)
como la diferencia incidencial \(120-90=30\).
En esta realización, la distribución sobre configuraciones microscópicas
es proporcional a \(\exp[-d(3X+4Y)]\).
Recibimos \(A_{\deg}=1000\alpha\),
\(C^*_{\deg}=2(\eta+\alpha)\) y ponemos
\(\xi=C^*_{\deg}/A_{\deg}=\tanh s\), en la rama
\(1/500<\xi<1\) de acción positiva. La rapidez de forma \(s\) es
adimensional; las dos coordenadas del cociente angular están expresadas
en grados.

\begin{proposition}[Curva de compatibilidad y lector de restitución]
\label{ins29:vii:curva}
Con \(d_*=50\pi f_\pi\alpha/243\), se cumple
\[
 d=d_*(1-\xi)=d_*(1-\tanh s),\qquad
 0<d<\frac{499}{500}d_*.
\]
La distribución anterior satisface, para \(m=1,\ldots,36\),
\[
 \frac{\Pr(M=m)}{\Pr(M=-m)}=e^{dm}.
\]
Por consiguiente, cada razón de probabilidades determina el mismo \(d\)
y restituye las coordenadas
\[
 \xi=1-\frac d{d_*},\qquad
 \eta=499\alpha-\frac{2430d}{\pi f_\pi},\qquad
 C^*_{\deg}=1000\alpha\left(1-\frac d{d_*}\right).
\]
\end{proposition}
\begin{proof}
La sustitución del par angular da
\(d=\pi f_\pi(499\alpha-\eta)/2430\).
Como \(\eta=500\alpha\xi-\alpha\), resulta la primera igualdad y,
por la rama declarada, el intervalo. Para la segunda, cada configuración
de ocupaciones \((X,Y)\) tiene una compañera \((Y,X)\) con la misma
multiplicidad. El cociente de sus pesos es
\(e^{-d(3X+4Y)+d(3Y+4X)}=e^{d(X-Y)}\).
Sumar sobre \(X-Y=m\) preserva ese factor. Aplicar el logaritmo y
despejar \(\xi\) y \(\eta\) proporciona las tres restituciones.
\end{proof}

Esta composición establece una comprobación conjunta: las treinta y seis
razones, divididas logarítmicamente por su respectivo \(m\), deben dar
la misma coordenada. La restitución actúa sobre salidas ya generadas y
mantiene su orden causal. Su resultado enlaza el sesgo de las ocupaciones
con la forma elíptica y la sección de acción, sin fijar un radio ni una
unidad energética adicional.

Los canales del funcional de Barbero se restituyen a continuación:
\[
 q_-=e^{-27d/f_\pi},\qquad
 q_+=D_A/q_-,\qquad D_A=e^{-100\pi\alpha/9}.
\]
En efecto, \(27d/f_\pi=\pi(A_{\deg}-C^*_{\deg})/180\), mientras
\(q_+q_-=D_A\). Sustituir estos canales y \(C^*_{\deg}(d)\) en
\eqref{md:compos:funcional-barbero} determina \(\gamma(d)\).
El coeficiente del área conserva entonces su composición
\(\gamma(d)a_0(d)L_\alpha^2\). Al recibir un autovalor de área
\(\mathcal A_n=\gamma a_0L_\alpha^2n\), el contraste de horizonte
da \(S_{\rm BH}(\mathcal A_n)/k_B=\gamma a_0n/4\).
La entropía del estado reducido conserva, además, sus probabilidades y
la bipartición. Ambas evaluaciones quedan disponibles con su objeto
respectivo: geometría del horizonte y distribución de entrelazamiento.
